\Askhsh{34422}{2}
\begin{ylh}{1}
    \item 3.1. Είδη και στοιχεία τριγώνων
    \item 3.2. 1ο Κριτήριο ισότητας τριγώνων
    \item 3.6. Κριτήρια ισότητας ορθογώνιων τριγώνων
    \item 3.11. Ανισοτικές σχέσεις πλευρών και γωνιών
    \item 4.6. Άθροισμα γωνιών τριγώνου
    \item 4.8. Άθροισμα γωνιών κυρτού ν-γώνου.
\end{ylh}
% ===================================================================
%  Αρχείο: 34422.tex (Fragment)
%  Αυτόματη μετατροπή μέσω doc2latex.py
% ===================================================================

ΘΕΜΑ 2

Δίνεται το ισοσκελές τρίγωνο ΑΒΓ με ΑΒ=ΑΓ. Φέρουμε, εκτός του τριγώνου, τις ημιευθείες Αx και Αy τέτοιες ώστε $Ax \perp AB$ και $Ay \perp AC$, όπως στο σχήμα που ακολουθεί. Οι κάθετες στην πλευρά ΒΓ στα σημεία Β και Γ τέμνουν τις Αx και Αy στα σημεία Δ και Ε αντίστοιχα.

\begin{alist}
\item Να αποδείξετε ότι ΒΔ=ΓΕ. \monades{12}

\item Αν η γωνία Β$\widehat{Α}$Γ είναι ίση με 80\textsuperscript{ο}, να υπολογίσετε τις γωνίες του τριγώνου που έχει για κορυφές τα σημεία Α, Ε και Δ. \monades{13}


\begin{center}
\begin{tikzpicture}[scale=1]
\tkzDefPoint(0,0){B}
\tkzDefPoint(4,0){C}
\tkzDefTriangle[two angles = 70 and 70](B,C) \tkzGetPoint{A}

% Ax perp AB
\tkzDefLine[orthogonal=through A](A,B) \tkzGetPoint{x_dir}
\tkzDefLine[orthogonal=through B](B,C) \tkzGetPoint{v_dir}
\tkzInterLL(A,x_dir)(B,v_dir) \tkzGetPoint{D}

% Ay perp AC
\tkzDefLine[orthogonal=through A](A,C) \tkzGetPoint{y_dir}
\tkzDefLine[orthogonal=through C](B,C) \tkzGetPoint{w_dir}
\tkzInterLL(A,y_dir)(C,w_dir) \tkzGetPoint{E}

\draw[l3] (A)--(B)--(C)--cycle;
\draw[l3] (A)--(D);
\draw[l3] (A)--(E);
\draw[l3] (B)--(D);
\draw[l3] (C)--(E);

\tkzDrawPoints(A,B,C,D,E)
\tkzLabelPoint[above](A){$A$}
\tkzLabelPoint[below](B){$B$}
\tkzLabelPoint[below](C){$\Gamma$}
\tkzLabelPoint[left](D){$\Delta$}
\tkzLabelPoint[right](E){$E$}

\tkzMarkRightAngle(B,A,D)
\tkzMarkRightAngle(E,A,C)
\tkzMarkRightAngle(C,B,D)
\tkzMarkRightAngle(B,C,E)
\end{tikzpicture}
\end{center}
\end{alist}
